\section{AMC 8 1997}
        \begin{enumerate}
        	\item (\textit{AMC 8 1997}) $\dfrac{1}{10} + \dfrac{9}{100} + \dfrac{9}{1000} + \dfrac{7}{10000} =$
$\text{(A)}\ 0.0026 \qquad \text{(B)}\ 0.0197 \qquad \text{(C)}\ 0.1997 \qquad \text{(D)}\ 0.26 \qquad \text{(E)}\ 1.997$

 \textbf{Answer: }C



\textbf{Solution}

Convert each fraction to a decimal.

 $\dfrac{1}{10} + \dfrac{9}{100} + \dfrac{9}{1000} + \dfrac{7}{10000}$
$0.1+0.09+0.009+0.0007=0.1997$
$\boxed{\textbf{(C)}}$



	\item (\textit{AMC 8 1997}) Ahn chooses a two-digit integer, subtracts it from 200, and doubles the result.  What is the largest number Ahn can get?


 $\text{(A)}\ 200 \qquad \text{(B)}\ 202 \qquad \text{(C)}\ 220 \qquad \text{(D)}\ 380 \qquad \text{(E)}\ 398$

 \textbf{Answer: }D



\textbf{Solution}

The smallest two-digit integer he can subtract from $200$ is $10$.  This will give the largest result for that first operation, and doubling it will keep it as the largest number possible.

 
 \[200-10=190\]
 \[190\times2=380\]
$\boxed{\textbf{(D)}}$



	\item (\textit{AMC 8 1997}) Which of the following numbers is the largest?


 $\text{(A)}\ 0.97 \qquad \text{(B)}\ 0.979 \qquad \text{(C)}\ 0.9709 \qquad \text{(D)}\ 0.907 \qquad \text{(E)}\ 0.9089$

 \textbf{Answer: }B



\textbf{Solution}

Using the process of elimination:

 Tenths digit: All tenths digits are equal, at $9$.

 Hundreths digit:$A$, $B$, and $C$ all have the same hundreths digit of $7$, and it is greater than the hundredths of $D$ or $E$ (which is $0$).  Eliminate both $D$ and $E$.

 Thousandths digit:$B$ has the largest thousandths digit of the remaining answers, and is the correct answer.  $A$ has an "invisible" thousandths digit of $0$, while $C$ also has a thousdandths digit of $0$.

 $\boxed{\textbf{(B)}}$



	\item (\textit{AMC 8 1997}) Julie is preparing a speech for her class.  Her speech must last between one-half hour and three-quarters of an hour.  The ideal rate of speech is 150 words per minute.  If Julie speaks at the ideal rate, which of the following number of words would be an appropriate length for her speech?


 $\text{(A)}\ 2250 \qquad \text{(B)}\ 3000 \qquad \text{(C)}\ 4200 \qquad \text{(D)}\ 4350 \qquad \text{(E)}\ 5650$

 \textbf{Answer: }E



\textbf{Solution}

Since there are $60$ minutes in an hour, we need a number of words between $(\frac{1}{2}\times 60)\times150=4500$ and $(\frac{3}{4}\times60)\times150=6750$.  The only such choice is $\boxed{\textbf{(E)}}$.

 


	\item (\textit{AMC 8 1997}) There are many two-digit multiples of 7, but only two of the multiples have a digit sum of 10.  The sum of these two multiples of 7 is


 $\text{(A)}\ 119 \qquad \text{(B)}\ 126 \qquad \text{(C)}\ 140 \qquad \text{(D)}\ 175 \qquad \text{(E)}\ 189$

 \textbf{Answer: }A



\textbf{Solution 1 (brute force)}

Writing out all two digit numbers that have a digital sum of $10$, you get $19, 28, 37, 46, 55, 64, 73, 82,$ and $91$.  The two numbers on that list that are divisible by $7$ are $28$ and $91$.  Their sum is $28+91=119$, choice $\boxed{\textbf{(A)}}$.

 



\textbf{Solution 2}

Writing out all the two digit multiples of $7$, you get $14, 21, 28, 35, 42, 49, 56, 63, 70, 77, 84,$ and $91$.  Again you find $28$ and $91$ have a digital sum of $10$, giving answer $\boxed{\textbf{(A)}}$.

 You may notice that adding $7$ either increases the digital sum by $7$, or decreases it by $2$, depending on whether there is carrying or not.

 


	\item (\textit{AMC 8 1997}) In the number $74982.1035$ the value of the \textit{place} occupied by the digit 9 is how many times as great as the value of the \textit{place} occupied by the digit 3?


 $\text{(A)}\ 1,000 \qquad \text{(B)}\ 10,000 \qquad \text{(C)}\ 100,000 \qquad \text{(D)}\ 1,000,000 \qquad \text{(E)}\ 10,000,000$

 \textbf{Answer: }C



\textbf{Solution 1}

The digit $9$ is $5$ places to the left of the digit $3$. Thus, it has a place value that is $10^5 = 100,000$ times greater.$\boxed{\textbf{(C)}}$




\textbf{Solution 2}

The digit $9$ is in the $100$s place.  The digit $3$ is in the $\frac{1}{1000}$ths place.  The ratio of these two numbers is $\frac{100}{\frac{1}{1000}} = 100\times1000 = 100,000$, and the answer is $\boxed{\textbf{(C)}}$



	\item (\textit{AMC 8 1997}) The area of the smallest square that will contain a circle of radius 4 is


 $\text{(A)}\ 8 \qquad \text{(B)}\ 16 \qquad \text{(C)}\ 32 \qquad \text{(D)}\ 64 \qquad \text{(E)}\ 128$

 \textbf{Answer: }D



\textbf{Solution}

Draw a square circumscribed around the circle.  (Alternately, the circle is inscribed in the square.)  If the circle has radius $4$, it has diameter $8$.  Two of the diameters of the circle will run parallel to the sides of the square.  Thus, the smallest square that contains it has side length $8$, and area $8\times8=64$.$\boxed{\textbf{(D)}}$



	\item (\textit{AMC 8 1997}) Walter gets up at 6:30 a.m., catches the school bus at 7:30 a.m., has 6 classes that last 50 minutes each, has 30 minutes for lunch, and has 2 hours additional time at school.  He takes the bus home and arrives at 4:00 p.m.  How many minutes has he spent on the bus?


 $\text{(A)}\ 30 \qquad \text{(B)}\ 60 \qquad \text{(C)}\ 75 \qquad \text{(D)}\ 90 \qquad \text{(E)}\ 120$

 \textbf{Answer: }B



\textbf{Solution}

There are $4\frac{1}{2}$ hours from 7:30 a.m. to noon, and $4$ hours from noon to 4 p.m., meaning Walter was away from home for a total of $8\frac{1}{2}$ hours, or $8.5\times 60 = 510$ minutes.

 Tallying up the times he was not on the bus, he has $6\times 50 = 300$ minutes in classes, $30$ minutes at lunch, and $2\times 60 = 120$ minutes of "additional time".  That is a total of $300 + 30 + 120 = 450$ minutes that Walter was away from home, but not on the bus.

 Therefore, Walter spent $510 - 450 = 60$ minutes on the bus, and the answer is $\boxed{B}$



	\item (\textit{AMC 8 1997}) Three students, with different names, line up in single file.  What is the probability that they are in alphabetical order from front-to-back?


 $\text{(A)}\ \dfrac{1}{12} \qquad \text{(B)}\ \dfrac{1}{9} \qquad \text{(C)}\ \dfrac{1}{6} \qquad \text{(D)}\ \dfrac{1}{3} \qquad \text{(E)}\ \dfrac{2}{3}$

 \textbf{Answer: }C



\textbf{Solution 1}

There are $3$ ways to pick the first student, $2$ ways to pick the second student, and $1$ way to pick the last student, for a total of $3! = 3\times2\times1 = 6$ ways to line the students up.

 Only $1$ of those ways is alphabetical.  Thus, the probability is $\frac{1}{6}$ or $\boxed{C}$




\textbf{Solution 2}

Give the students uncreative names like Abby Adams, Bob Bott, and Carol Crock:  $A$, $B$, and $C$, replacing the alphabetically first student's name with $A$, and the alphabetically last student's name with $C$.  List all the ways they can line up:

 $ABC$
$ACB$
$BAC$
$BCA$
$CAB$
$CBA$

 Only $1$ of those $6$ lists is in alphabetical order, giving an answer of $\frac{1}{6}$ or $\boxed{C}$




\textbf{Solution 3}

We look at the first place in the line, there is only one correct person who belongs in this place, but without the restriction there can be three different people who can go in this spot. So we have a $\frac{1}{3}$ chance of getting the first place right. Now look at the second place, there is only one correct person who belongs in this place, but without the restriction there can be two different people who can go in this place, as one person is already in the first place. So we have a $\frac{1}{2}$ chance of getting the second place right.Now look at the third place, there is only one correct person who belongs in this place, but since all the other spots are filled there is only one person who can be in this place at all. So we have a $\frac{1}{1}$ chance of getting the third place right.Since the total probability of getting the order depends on the probability of each place having the right person we should multiply these three fractions.$\frac{1}{3}\times\frac{1}{2}\times\frac{1}{1}$ is the probability of getting the order right.Simplifying this we get that the answer is $\frac{1}{6}$ or $\boxed{C}$.~blankbox

 


	\item (\textit{AMC 8 1997}) What fraction of this square region is shaded?  Stripes are equal in width, and the figure is drawn to scale.


 
                    \begin{figure}[H]
                        \centering
                        \includegraphics[width=0.5\textwidth]{Problem/AMC_8_AJHSME_1997_10_Problem_Figure_1.png}
                    \end{figure}
                    
$\text{(A)}\ \dfrac{5}{12} \qquad \text{(B)}\ \dfrac{1}{2} \qquad \text{(C)}\ \dfrac{7}{12} \qquad \text{(D)}\ \dfrac{2}{3} \qquad \text{(E)}\ \dfrac{5}{6}$

 \textbf{Answer: }C



\textbf{Solution 1}

Fill in grid lines to make the shape look like a small 6x6 checkerboard.

 There are $6\times6 = 36$ squares total.  The smallest shaded region has $3$ squares.  The medium shaded region has $7$ squares.  The largest shaded region has $11$ squares.  The total shaded area is $3 + 7 + 11 = 21$ squares.

 Thus, $\frac{21}{36} = \frac{7}{12}$ of the square is shaded, and the answer is $\boxed{C}$




\textbf{Solution 2}

Similar to the first solution, fill in the gridlines to make $36$ squares total.

 Instead of counting shaded squares, count unshaded squares.  There are $1 + 5 + 9 = 15$ unshaded squares.

 Thus $\frac{15}{36} = \frac{5}{12}$ of the square is unshaded, and $1 - \frac{5}{12} = \frac{7}{12}$ of the square is shaded.

 


	\item (\textit{AMC 8 1997}) Define $\boxed{N}$ as the number of whole number divisors of $N$.  For example, $\boxed{3}=2$ because 3 has two divisors, 1 and 3.  Find the value of


 
 \[\boxed{\boxed{11}\times\boxed{20}}.\]
$\text{(A)}\ 6 \qquad \text{(B)}\ 8 \qquad \text{(C)}\ 12 \qquad \text{(D)}\ 16 \qquad \text{(E)}\ 24$

 \textbf{Answer: }A



\textbf{Solution}

$\boxed{\boxed{11}\times\boxed{20}}$
$11$ has $2$ factors, and $20$ has $\{1, 2, 4, 5, 10, 20\}$ as factors, for a total of $6$ factors.  Plugging $\boxed{11} = 2$ and $\boxed {20}= 6$:

 $\boxed{2 \times 6}$
$\boxed{12}$
$12$ has factors of $\{1, 2, 3, 4, 6, 12\}$, so $\boxed{12} = 6$, and the answer is $\boxed{A}$



	\item (\textit{AMC 8 1997}) $\angle 1 + \angle 2 = 180^\circ$
$\angle 3 = \angle 4$


 Find $\angle 4.$

                    \begin{figure}[H]
                        \centering
                        \includegraphics[width=0.5\textwidth]{Problem/AMC_8_AJHSME_1997_12_Problem_Figure_1.png}
                    \end{figure}
                    
$\text{(A)}\ 20^\circ \qquad \text{(B)}\ 25^\circ \qquad \text{(C)}\ 30^\circ \qquad \text{(D)}\ 35^\circ \qquad \text{(E)}\ 40^\circ$

 \textbf{Answer: }D



\textbf{Solution}

Using the left triangle, we have:

 $\angle 1 + 70 + 40 = 180$
$\angle 1 = 180 - 110$
$\angle 1 = 70$

 Using the given fact that $\angle 1 + \angle 2 = 180$, we have $\angle 2 = 180 - 70 = 110$.

 Finally, using the right triangle, and the fact that $\angle 3 = \angle 4$, we have:

 $\angle 2 + \angle 3 + \angle 4 = 180$
$110 + \angle 3 + \angle 4 = 180$
$110 + \angle 4+ \angle 4 = 180$
$2\angle 4 = 70$
$\angle 4 = 35$

 Thus, the answer is $\boxed{D}$

 \textbf{Side Note to Make Solution Faster} You can use Exterior Angle Theorem to directly say that $\angle 1 = \angle 3+\angle 4$ instead of needing to use up time by doing $\angle 2 + \angle 3 + \angle 4 = 180$, so on. Using exterior angle theorem goes straight to $\angle 1 = \angle 3+\angle 4$, and therefore, $2(\angle 4) = 70^{\circ}$, as we know $180^{\circ}-40^{\circ}-70^{\circ} = \angle 1 = 70^{\circ}$, and $\angle 3 = \angle 4$. This means that $\angle 4 = 35^{\circ}$, so the answer is  $\boxed{D}$.

 


	\item (\textit{AMC 8 1997}) Three bags of jelly beans contain 26, 28, and 30 beans.  The ratios of yellow beans to all beans in each of these bags are $50\%$, $25\%$, and $20\%$, respectively.  All three bags of candy are dumped into one bowl.  Which of the following is closest to the ratio of yellow jelly beans to all beans in the bowl?


 $\text{(A)}\ 31\% \qquad \text{(B)}\ 32\% \qquad \text{(C)}\ 33\% \qquad \text{(D)}\ 35\% \qquad \text{(E)}\ 95\%$

 \textbf{Answer: }A



\textbf{Solution}

In bag $A$, there are $26$ jellybeans, and $50\%$ are yellow.  That means there are $26\times 50\% = 26\times 0.50 = 13$ yellow jelly beans in this bag.

 In bag $B$, there are $28$ jellybeans, and $25\%$ are yellow.  That means there are $28\times 25\% = 28\times 0.25 = 7$ yellow jelly beans in this bag.

 In bag $C$, there are $30$ jellybeans, and $20\%$ are yellow.  That means there are $30\times 20\% = 30\times 0.20 = 6$ yellow jelly beans in this bag.

 In all three bags, there are $13 + 7 + 6 = 26$ yellow jelly beans in total , and $26 + 28 + 30 = 84$ jelly beans of all types in total.

 Thus, $\frac{26}{84} = \frac{26}{84}\cdot 100\% = \frac{13 }{ 42} \cdot 100\% = 30.9\%$ of all jellybeans are yellow.  Thus, the correct answer is $\boxed{A}$



	\item (\textit{AMC 8 1997}) There is a set of five positive integers whose average (mean) is 5, whose median is 5, and whose only mode is 8.  What is the difference between the largest and smallest integers in the set?


 $\text{(A)}\ 3 \qquad \text{(B)}\ 5 \qquad \text{(C)}\ 6 \qquad \text{(D)}\ 7 \qquad \text{(E)}\ 8$

 \textbf{Answer: }D



\textbf{Solution}

Let $a$ and $b$ be the two smallest integers, with $a \leq b$.  Since $5$ is the median and $8$ is the unique mode, the integers, from smallest to largest, are $a, b, 5, 8, 8$.  Since the mean is $5$, the sum of all five integers is $5 \cdot 5 = 25$, and thus $a + b = 4$.  Since $8$ is the unique mode, $a \neq b$.  So $a = 1$ and $b = 3$, and thus the range is $8-1=\boxed{\textbf{(D)}\ 7}$.

 


	\item (\textit{AMC 8 1997}) Each side of the large square in the figure is trisected (divided into three equal parts).  The corners of an inscribed square are at these trisection points, as shown.  The ratio of the area of the inscribed square to the area of the large square is


 
                    \begin{figure}[H]
                        \centering
                        \includegraphics[width=0.5\textwidth]{Problem/AMC_8_AJHSME_1997_15_Problem_Figure_1.png}
                    \end{figure}
                    
$\text{(A)}\ \dfrac{\sqrt{3}}{3} \qquad \text{(B)}\ \dfrac{5}{9} \qquad \text{(C)}\ \dfrac{2}{3} \qquad \text{(D)}\ \dfrac{\sqrt{5}}{3} \qquad \text{(E)}\ \dfrac{7}{9}$

 \textbf{Answer: }B



\textbf{Solution 1}

Define $1$ unit as the distance from vertex of the large square to the closest trisection point.  Then the large square has side length $3$, so its area is $9$.  

 Each of the four right triangles shown has area $\frac{1}{2}\cdot 2\cdot 1 = 1$, so the area of the small square is $9 - 4\cdot 1 = 5$.  The ratio of the areas of the two squares is thus $\boxed{\textbf{(B)}\ \frac{5}{9}}$.

 



\textbf{Solution 2}

Let the side of the big square be $3x$.  The area of this square is $3x\cdot 3x = 9x^2$

 You can use the Pythagorean Theorem on the any of the right triangles to find the length of the side of the inscribed square.  One leg is $x$, while the other leg is $2x$.  Thus, the length of the hypotenuse $c$, which is also the side of the inscrubed square, is:

 $a^2 + b^2 = c^2$
$x^2 + (2x)^2 = c^2$
$x^2 + 4x^2 = c^2$
$5x^2 = c^2$

 Once you get $c^2$, you can stop, because $c^2$ is the area of the inscribed square.  The ratio of the areas  of the small square to the big square is $\frac{5x^2}{9x^2} = \frac{5}{9}$, giving option $\boxed{B}$.

 



\textbf{Solution 3}

Fill in the gridlines of the diagram.

 
                    \begin{figure}[H]
                        \centering
                        \includegraphics[width=0.5\textwidth]{Solution/AMC_8_AJHSME_1997_15_Solution_Figure_1.png}
                    \end{figure}
                    

 Let each grid space have length $1$. The inscribed square consists of four congruent triangles and one square.

 The triangles each have area $\frac{1}{2} \cdot 2\cdot 1=1$, and in total they cover $4 \cdot 1 = 4$ square units of area. The central square's area is $1 \cdot 1=1$. Therefore, the area of the inscribed square is $4+1=5$.

 The big square contains $3 \cdot 3=9$ square units of area. Thus, the ratio is equal to $\frac{5}{9}$, or $\boxed{B}$.

 


 


	\item (\textit{AMC 8 1997}) Penni Precisely buys $100 worth of stock in each of three companies: Alabama Almonds, Boston Beans, and California Cauliflower.  After one year, AA was up 20%, BB was down 25%, and CC was unchanged.  For the second year, AA was down 20% from the previous year, BB was up 25% from the previous year, and CC was unchanged.  If A, B, and C are the final values of the stock, then


 $\text{(A)}\ A=B=C \qquad \text{(B)}\ A=B<C \qquad \text{(C)}\ C<B=A$
$\text{(D)}\ A<B<C \qquad \text{(E)}\ B<A<C$

 \textbf{Answer: }E



\textbf{Solution 1}

AA is $100 \cdot 120\% = 100\cdot 1.2 = 120$ after one year.  After the second year, AA is $120\cdot 0.8 = 96$.

 BB is $100 \cdot 75\% = 100\cdot 0.75 = 75$ after one year.  After the second year, BB is $75\cdot 1.25 = 93.75$

 CC remains unchanged throughout, and stays at $100$.

 Thus, $B < A < C$, and the right answer is $\boxed{E}$




\textbf{Solution 2}

AA will be $100 \cdot \frac{6}{5} \cdot \frac{4}{5} = 100\cdot \frac{24}{25}$ at the end.

 BB will be $100 \cdot \frac{3}{4} \cdot \frac{5}{4} = 100 \cdot \frac{15}{16}$ at the end.

 CC will be unchanged at $100$.

 Since all the fractions are under $1$, $C$ will be highest value.

 Since $\frac{24}{25}$ is only $\frac{1}{25}$ away from $1$, while $\frac{15}{16}$ is $\frac{1}{16}$ away from $1$, $\frac{24}{25}$ is closer to $1$, and will be closer to the original value.

 Thus, $B < A < C$, and the right answer is $\boxed{E}$



	\item (\textit{AMC 8 1997}) A cube has eight vertices (corners) and twelve edges.  A segment, such as $x$, which joins two vertices not joined by an edge is called a diagonal.  Segment $y$ is also a diagonal.  How many diagonals does a cube have?


 
                    \begin{figure}[H]
                        \centering
                        \includegraphics[width=0.5\textwidth]{Problem/AMC_8_AJHSME_1997_17_Problem_Figure_1.png}
                    \end{figure}
                    
$\text{(A)}\ 6 \qquad \text{(B)}\ 8 \qquad \text{(C)}\ 12 \qquad \text{(D)}\ 14 \qquad \text{(E)}\ 16$

 \textbf{Answer: }E



\textbf{Solution 1}

On each face, there are $2$ diagonals like $x$.  There are $6$ faces on a cube.  Thus, there are $2\times 6 = 12$ diagonals that are "x-like".

 Every "y-like" diagonal must connect the bottom of the cube to the top of the cube.  Thus, for each of the $4$ bottom vertices of the cube, there is a different "y-like" diagonal.  So there are $4$ "y-like" diagonals.

 This gives a total of $12 + 4 = 16$ diagonals on the cube, which is answer $\boxed{E}$.

 



\textbf{Solution 2}

There are $8$ vertices on a cube(A, B, C, D, E, F, G, H).  If you pick one vertice, it will have 7 line segments associated with it, so you will have $\frac{8\cdot7}{2} = 28$ segments within the cube.  The division by $2$ is necessary because you counted both the segment from $A$ to $B$ and the segment from $B$ to $A$.

 But not all $28$ of these segments are diagonals.  Some are edges.  There are $4$ edges on the top, $4$ edges on the bottom, and $4$ edges that connect the top to the bottom.  So there are $12$ edges total, meaning that there are $28 - 12 = 16$ segments that are not edges.  All of these segments are diagonals, and thus the answer is $\boxed{E}$.

 



\textbf{Solution 3}

Consider picking one point on the corner of the cube.  That point has $3$ "x-like" diagonals , and $1$ "y-like" diagonal that ends on the opposite vertex.  Thus, each vertex has $3 + 1 = 4$ diagonals associated with it.  There are $8$ vertices on the cube, giving a total of $8 \cdot 4 = 32$ diagonals.

 However, each diagonal was counted as both a "starting point" and an "ending point".  So there are really $\frac{32}{2} = 16$ diagonals, giving an answer of $\boxed{E}$.

 


	\item (\textit{AMC 8 1997}) At the grocery store last week, small boxes of facial tissue were priced at 4 boxes for $5.  This week they are on sale at 5 boxes for $4.  The percent decrease in the price per box during the sale was closest to


 $\text{(A)}\ 30\% \qquad \text{(B)}\ 35\% \qquad \text{(C)}\ 40\% \qquad \text{(D)}\ 45\% \qquad \text{(E)}\ 65\%$

 \textbf{Answer: }B



\textbf{Solution}

Last week, each box was $\frac{5}{4} = 1.25$

 This week, each box is $\frac{4}{5} = 0.80$

 Percent decrease is given by $\frac{X_{old} - X_{new}}{X_{old}} \cdot 100\%$

 This, the percent decrease is $\frac{1.25 - 0.8}{1.25}\cdot 100\% = \frac{45}{125} \cdot 100\% = 36\%$, which is closest to $\boxed{B}$



	\item (\textit{AMC 8 1997}) If the product $\dfrac{3}{2}\cdot \dfrac{4}{3}\cdot \dfrac{5}{4}\cdot \dfrac{6}{5}\cdot \ldots\cdot \dfrac{a}{b} = 9$, what is the sum of $a$ and $b$?


 $\text{(A)}\ 11 \qquad \text{(B)}\ 13 \qquad \text{(C)}\ 17 \qquad \text{(D)}\ 35 \qquad \text{(E)}\ 37$

 \textbf{Answer: }D



\textbf{Solution 1}

Notice that the numerator of the first fraction cancels out the denominator of the second fraction, and the numerator of the second fraction cancels out the denominator of the third fraction, and so on.

 The only numbers left will be $a$ in the numerator from the last fraction and $2$ in the denominator from the first fraction.  (The $b$ will cancel with the numerator of the preceeding number.)  Thus, $\frac{a}{2} = 9$, and $a=18$.

 Since the numerator is always one more than the denominator, $b = a-1 = 17$, and $a+ b = 18 + 17 = 35$, giving an answer of $\boxed{D}$




\textbf{Solution 2}

Find a pattern.  If $a=4$, then the expression is just the first two terms, which is $2$.

 If $a=5$, then the expression is the first three terms, giving $2.5$.

 If $a=6$, the expression is the first four terms, giving $3$.

 If $a=7$, the expression will be the first five terms, giving $3.5$.

 Conjecture that the expression is always going to equal $\frac{a}{2}$, and thus when $a=18$, the expression will be $9$, as desired.

 As above, when $a=18$, $b=17$, and the sum is $35$, or $\boxed{D}$



	\item (\textit{AMC 8 1997}) A pair of 8-sided dice have sides numbered 1 through 8.  Each side has the same probability (chance) of landing face up.  The probability that the product of the two numbers that land face-up exceeds 36 is


 $\text{(A)}\ \dfrac{5}{32} \qquad \text{(B)}\ \dfrac{11}{64} \qquad \text{(C)}\ \dfrac{3}{16} \qquad \text{(D)}\ \dfrac{1}{4} \qquad \text{(E)}\ \dfrac{1}{2}$

 \textbf{Answer: }A



\textbf{Solution}

There are $8 \times 8 = 64$ combinations to examine.

 If one die is $1$, then even with an $8$ on the other die, no combinations will work.

 If one die is $2$, then even with an $8$ on the other die, no combinations will work.

 If one die is $3$, then even with an $8$ on the other die, no combinations will work.

 If one die is $4$, then even with an $8$ on the other die, no combinations will work.

 If one die is $5$, then the other die must be an $8$ to have a product over $36$.  Thus, $(5,8)$ works.

 If one die is $6$, then the other die must be either $7$ or $8$ to have a product over $36$.  Thus, $(6, 7)$ and $(6, 8)$ both work.

 If one die is $7$, then the other die can be $6, 7,$ or $8$ to have a product over $36$.  Thus, $(7, 6)$, $(7, 7)$, and $(7, 8)$ all work.

 If one die is $8$, then the other die can be $5, 6, 7,$ or $8$ to have a product over $36$.  Thus, $(8, 5), (8,6), (8,7),$ and $(8,8)$ work.

 There are a total of $1 + 2 + 3 + 4 = 10$ combinations that work out of a total of $64$ possibilities.

 Thus, the answer is $\frac{10}{64} = \frac{5}{32}$, and the answer is $\boxed{A}$

 But  this  may  be  slightly  complicated,  so  you  can  also  do  solution  2 (scroll to bottom).

 



\textbf{Solution 2}

Make a chart with all products. 10 work out of 64. Simplify for 5/32 or $\boxed{A}$.

 


	\item (\textit{AMC 8 1997}) Each corner cube is removed from this $3\text{ cm}\times 3\text{ cm}\times 3\text{ cm}$ cube.  The surface area of the remaining figure is


 
                    \begin{figure}[H]
                        \centering
                        \includegraphics[width=0.5\textwidth]{Problem/AMC_8_AJHSME_1997_21_Problem_Figure_1.png}
                    \end{figure}
                    
$\text{(A)}\ 19\text{ sq.cm} \qquad \text{(B)}\ 24\text{ sq.cm} \qquad \text{(C)}\ 30\text{ sq.cm} \qquad \text{(D)}\ 54\text{ sq.cm} \qquad \text{(E)}\ 72\text{ sq.cm}$

 \textbf{Answer: }D



\textbf{Solution}

The original cube has $6$ square surfaces that each have an area of $3^2 = 9$, for a toal surface area of $6\cdot 9 = 54$.

 Since no two corner cubes touch, we can examine the effect of removing each corner cube individually.

 Each corner cube contribues $3$ faces each of surface area $1$ to the big cube, so the surface area is decreased by $3$ when the cube is removed.

 However, when the cube is removed, $3$ faces on the 3x3x3 cube will be revealed, increasing the surface area by $3$.

 Thus, the surface area does not change with the removal of a corner cube, and it remains $54$, which is answer $\boxed{D}$.

 


 


	\item (\textit{AMC 8 1997}) A two-inch cube $(2\times 2\times 2)$ of silver weighs 3 pounds and is worth $200.  How much is a three-inch cube of silver worth?


 $\text{(A)}\ 300\text{ dollars} \qquad \text{(B)}\ 375\text{ dollars} \qquad \text{(C)}\ 450\text{ dollars} \qquad \text{(D)}\ 560\text{ dollars} \qquad \text{(E)}\ 675\text{ dollars}$

 \textbf{Answer: }E



\textbf{Solution}

The two-inch cube has a volume of $8$ cubic inches, and the three-inch cube has a volume of $27$ cubic inches. Thus, the three-inch cube has a weight that is $\frac{27}{8}$ times that of the two-inch cube. Then its value is $\frac{27}{8} \cdot 200 = \boxed{\textbf{(E) }\text{675 dollars}}$.

 


	\item (\textit{AMC 8 1997}) There are positive integers that have these properties:


 
            \begin{itemize}
        \item the sum of the squares of their digits is 50, and

        \item each digit is larger than the one to its left.
            \end{itemize}
            The product of the digits of the largest integer with both properties is


 $\text{(A)}\ 7 \qquad \text{(B)}\ 25 \qquad \text{(C)}\ 36 \qquad \text{(D)}\ 48 \qquad \text{(E)}\ 60$

 \textbf{Answer: }C



\textbf{Solution}

Five-digit numbers will have a minimum of $1^2 + 2^2 + 3^2 + 4^2 + 5^2 = 55$ as the sum of their squares if the five digits are distinct and non-zero.  If there is a zero, it will be forced to the left by rule \#2.

 No digit will be greater than $7$, as $8^2 = 64$.

 Trying four digit numbers $WXYZ$, we have $w^2 + x^2 + y^2 + z^2 = 50$ with $0 < w < x < y < z < 8$
$z=7$ will not work, since the other digits must be at least $1^2 + 2^2 + 3^2 = 14$, and the sum of the squares would be over $50$.

 $z=6$ will give $w^2 + x^2 + y^2 = 14$.  $(w,x,y) = (1,2,3)$ will work, giving the number $1236$.  No other number with $z=6$ will work, as $w, x,$ and $y$ would have to be greater. 

 $z=5$ will give $w^2 + x^2 + y^2 = 25$.  $y=4$ forces $x=3$ and $w=0$, which has a leading zero, and then we have 345 which is a 3-digit number. 

 $z=4$ can only give the number $1234$, which does not satisfy the condition of the problem.

 Thus, the number in question is $1236$, and the product of the digits is $36$, giving $\boxed{C}$ as the answer.

 


	\item (\textit{AMC 8 1997}) Diameter $ACE$ is divided at $C$ in the ratio $2:3$.  The two semicircles, $ABC$ and $CDE$, divide the circular region into an upper (shaded) region and a lower region.  The ratio of the area of the upper region to that of the lower region is


 
                    \begin{figure}[H]
                        \centering
                        \includegraphics[width=0.5\textwidth]{Problem/AMC_8_AJHSME_1997_24_Problem_Figure_1.png}
                    \end{figure}
                    
$\text{(A)}\ 2:3 \qquad \text{(B)}\ 1:1 \qquad \text{(C)}\ 3:2 \qquad \text{(D)}\ 9:4 \qquad \text{(E)}\ 5:2$

 \textbf{Answer: }C



\textbf{Solution}


                    \begin{figure}[H]
                        \centering
                        \includegraphics[width=0.5\textwidth]{Solution/AMC_8_AJHSME_1997_24_Solution_Figure_1.png}
                    \end{figure}
                    

 Draw $\overline{AE}$ to divide the big circle in half.  Assign $AC = 4$ and $CE = 6$ so that the radii work out to be integers.  ($4x$ and $6x$ can be used instead, but the $x$ will cancel in the ratio.)  

 The shaded region is equal to the area of semicircle $\overarc{AE}$ on top, plus the area of the semicircle $CDE$ on the bottom, minus the area of semicircle $\overarc{ABC}$ on top.

 The radii of those three semicircles are $5, 3,$ and $2$, respectively.Thus, the area of the shaded region is $\frac{1}{2}\pi\cdot5^2 + \frac{1}{2}\pi\cdot3^2 - \frac{1}{2}\pi\cdot2^2 =\frac{\pi}{2}(5^2 + 3^2 - 2^2)=15\pi$

 The total area of the circle is $\pi \cdot 5^2 = 25\pi$.Thus, the unshaded area is $25\pi - 15\pi = 10\pi$.Therefore the ratio of shaded:unshaded is $15\pi : 10\pi =\boxed{ \text{(C)}\ 3:2}$.

 


	\item (\textit{AMC 8 1997}) All of the even numbers from 2 to 98 inclusive, excluding those ending in 0, are multiplied together.  What is the rightmost digit (the units digit) of the product?


 $\text{(A)}\ 0 \qquad \text{(B)}\ 2 \qquad \text{(C)}\ 4 \qquad \text{(D)}\ 6 \qquad \text{(E)}\ 8$

 \textbf{Answer: }D



\textbf{Solution 1}

All the tens digits of the product will be irrelevant to finding the units digit.  Thus, we are searching for the units digit of $(2\cdot 4\cdot 6 \cdot 8) \cdot (2 \cdot 4 \cdot 6 \cdot 8) \cdot (2\cdot 4\cdot 6 \cdot 8) \cdot ...$

 There will be $10$ groups of $4$ numbers.  The number now can be rewritten as $(2\cdot 4 \cdot 6 \cdot 8)^{10}$.

 Multiplying, we get $384^{10}$.

 Again, we can disregard the tens and hundreds digit of $384$, since we only want the units digit of the number, leaving $4^{10}$.

 Now, we try to find a pattern to the units digit of $4^n$.  To compute this quickly, we once again discard all tens digits and higher.

 $4^1 = 4$.

 $4^2 = 4\cdot 4 = 1\underline{6}$, discard the $1$.

 $4^3 = 1 \cdot 4 = \underline{4}$
$4^4 = 4 \cdot 4 = 1\underline{6}$, discard the $1$.

 $4^5 = 1 \cdot 4 = \underline{4}$

 Those equalities are, in reality, congruences $\mod {10}$.

 Thus, the pattern of the units digits is $\{4, 6, 4, 6, 4, 6, 4, 6, 4, 6\}$.  The cycle repeats so that term $n$ is the same as term $n+2$.  The tenth number in the cycle is $6$, giving an answer of $\boxed{D}$.

 



\textbf{Solution 2}

Again, the value we seek is equal to $(2\cdot 4\cdot 6\cdot 8)^{10}\mod 10$.We can use equivalence to simplify.

 $(2\cdot 4\cdot 6\cdot 8)^{10}\mod 10$
$\equiv (2\cdot 4\cdot (-4)\cdot (-2))^{10} \mod 10$
$\equiv 64^{10} \mod 10$
$\equiv 4^{10} \mod 10$
$\equiv 2^{20} \mod 10$
$\equiv 8^6\cdot 2^2 \mod 10$
$\equiv (-2)^6\cdot 2^2 \mod 10$
$\equiv 2^8 \mod 10$
$\equiv 256 \mod 10$
$\equiv 6 \mod 10$

 Thus our answer is $\boxed{D}$.

 


\end{enumerate}